\pdfoutput=1
\documentclass[11pt]{article}

\usepackage{ACL2023}
\usepackage{times}
\usepackage{latexsym}
\usepackage[T1]{fontenc}

\usepackage[utf8]{inputenc}
\usepackage{microtype}
\usepackage{inconsolata}
\usepackage{fancyhdr}
\usepackage{natbib}
\usepackage{booktabs}
\usepackage{array}

\pagestyle{fancy}
\fancyhf{}
\fancyhead[LE,RO]{EPFL -- Spring Semester 2026}
\fancyhead[RE,LO]{Modern NLP -- Open Project Final Report}
\rfoot{Page \thepage \hspace{1pt}}
\renewcommand{\headrulewidth}{1pt}

%%%%%%%%%%%%%%%%%%% TITLE AND AUTHORS %%%%%%%%%%%%%%%%%%%

\title{Faithful RAG: Evaluating Citation Accuracy and Retrieval Robustness in Domain-Specific Scientific Literature}

\author{\normalfont
Andrea Trugenberger | 357615 | \texttt{andrea.trugenberger@epfl.ch} \\
Elie Bruno | 355932 | \texttt{elie.bruno@epfl.ch} \\
Faruk Zahiragi\'c | 415360 | \texttt{faruk.zahiragic@epfl.ch} \\
Yusif Askari | 413862 | \texttt{yusif.askari@epfl.ch} \\
Team CiteRight \\
}

%%%%%%%%%%%%%%%%%%% FINAL REPORT %%%%%%%%%%%%%%%%%%%

\begin{document}
\maketitle
\thispagestyle{fancy}

\begin{abstract}
% TODO (group): <300 words. Motivate the faithfulness gap in RAG, summarize
% the four-component pipeline (retrieval ablation, NLI/judge faithfulness,
% chunked-RAG vs long-context, CRAG router), and headline two numbers:
% (i) E5-large coarse+rerank as new SOTA (0.878 MRR, +8.6pp over OpenAI baseline),
% (ii) 55--61pp NLI-vs-judge faithfulness disagreement showing NLI alone fails
% as a faithfulness gate.
Faithfulness in retrieval-augmented generation (RAG) is typically measured by
end-to-end answer-level metrics that conflate retrieval quality with grounded
generation. We disentangle the two on a 93-claim, character-span--validated
gold benchmark over intellectual-property and innovation-policy papers, and
report on four interventions: a sixteen-configuration retrieval ablation, a
faithfulness verifier combining DeBERTa-v3 NLI with an LLM judge, a chunked-RAG
vs.\ 1M-token long-context comparison with cost normalisation, and a Corrective
RAG router \citep{yan2024corrective}. E5-large coarse+rerank is the new
corpus-level SOTA at 0.878 MRR (+8.6\,pp over a closed-source OpenAI baseline),
the NLI and LLM-judge heads disagree by 55--61\,pp exposing NLI as an
insufficient gate, and CRAG lifts retrieval correctness from 60\% to 80\% on
a mixed-quality slice. % TODO: finalize once human-eval and chunker-ablation numbers land.
\end{abstract}


\section{Introduction}
\label{sec:intro}
% TODO (group): per TA feedback, this section must MOTIVATE the work before
% diving into approach. Recommended structure:
%   1. The problem: RAG systems are deployed under the assumption that
%      retrieved context grounds the answer, but recent work shows substantial
%      fabrication even when retrieval succeeds.
%   2. Why current methods fall short: end-to-end metrics (EM, BLEU, judge
%      scores) conflate retrieval recall with generative grounding; NLI-based
%      faithfulness checks suffer known failure modes on long-form
%      multi-claim outputs.
%   3. Our approach: we build a character-span-grounded gold benchmark
%      (93 QA pairs, 5 difficulty buckets, 6 adversarial controls) over
%      domain-specific scientific literature, and run four orthogonal
%      interventions that each isolate one failure axis: retrieval, judging,
%      long-context tradeoff, and corrective routing.
%   4. Contributions (bullet form): (i) gold benchmark, (ii) retrieval-config
%      SOTA, (iii) NLI-vs-judge gap quantification, (iv) CRAG ablation,
%      (v) human eval validating the judge.
%   5. Key results preview (one sentence each, mirroring the abstract).


\section{Approach}
\label{sec:approach}
% TODO (group): describe the four-component pipeline. Per intermediate
% feedback ("a brief note on how the experiments are planned across the team,
% and what the INCORRECT-branch fallback will actually do"), this section
% must:
%   - name per-member ownership of each component (retrieval=Elie, NLI/judge
%     =Andrea, CRAG=Faruk, RAGAS/long-context=Yusif)
%   - spell out the CRAG INCORRECT-branch fallback action (web search?
%     abstain? rewrite?) — currently underspecified.

\subsection{Gold benchmark}
\label{sec:approach:gold}
% TODO: describe gold_qa.json schema: 93 claim-level QA pairs across 5
% difficulty buckets (single-hop, multi-hop, unanswerable, IAA subset,
% adversarial controls q900-q905), with byte-validated supporting_spans
% (char_start/char_end into LF-normalised document.md). Note IAA subset
% targets Cohen-kappa across reviewers.

\subsection{Retrieval pipeline}
\label{sec:approach:retrieval}
% TODO: brief description of HybridRetriever code path; embedder families;
% optional ZeroEntropy cross-encoder reranker. Defer numbers to Results.

\subsection{Faithfulness verification}
\label{sec:approach:faith}
% TODO: NLI head (DeBERTa-v3-large-mnli) + LLM-judge head; how the two are
% combined; thresholds.

\subsection{Corrective RAG router}
\label{sec:approach:crag}
% TODO: CORRECT / AMBIGUOUS / INCORRECT branch logic; thresholds (tau_h, tau_l);
% the INCORRECT-branch fallback (specify per intermediate feedback).


\section{Experiments}
\label{sec:exp}

\subsection{Data}
\label{sec:exp:data}
% TODO: dataset sources (intellectual-property + innovation-policy papers
% sourced from arXiv / OpenAlex), preprocessing (LF normalisation, document.md
% extraction), annotation protocol (annotator + reviewer + IAA subset),
% adversarial controls rationale, licensing / ethical considerations on
% the source corpus.

\subsection{Evaluation}
\label{sec:exp:eval}
% TODO: per-component metrics:
%   - Retrieval: hit@k, MRR, nDCG (paper-level), open-recall.
%   - Faithfulness: NLI label vs LLM-judge label vs human label (kappa).
%   - RAGAS: faithfulness, answer_relevancy, context_precision, context_recall.
%   - CRAG: pre/post correct-retrieval rate; latency.
%   - Cost: USD per query, normalized accuracy-per-dollar.
% Baselines: OpenAI text-embedding-3-small (closed-source), BGE-M3, ColBERTv2;
% long-context Gemini 2.5 Flash; closed-book GPT-4o-mini / Claude-3.5-haiku /
% DeepSeek-chat.

\subsection{Experimental details}
\label{sec:exp:details}
% TODO: hardware (RCP cluster, GPU types), training/inference time per
% experiment, prompt versions, library versions (pin uv.lock), random seeds.


\section{Results and Analysis}
\label{sec:results}

\subsection{Retrieval ablation}
\label{sec:retrieval}

The 93-claim gold set (5 categories incl.\ multi-hop, byte-validated character spans, 6 adversarial controls q900--q905 that break the single-class Cohen-$\kappa$ paradox \citep{cohen1960kappa} from the first IAA round) drives a sixteen-configuration ablation. Configurations cross four embedder families (OpenAI \texttt{text-embedding-3-small}; BGE-M3 \citep{chen2024bgem3}; E5-large \citep{wang2024e5}; ColBERTv2 \citep{santhanam2022colbertv2}), two chunk granularities (coarse $\sim$2000\,chars, fine $\sim$300\,chars), and an optional ZeroEntropy cross-encoder reranker, sharing one \texttt{HybridRetriever} code path that reports hit@$k$, MRR, and nDCG.

\begin{table}[h]
\centering
\footnotesize
\begin{tabular}{lccc}
\toprule
config (\texttt{coarse\_rerank}) & hit@5 & hit@10 & MRR \\
\midrule
OpenAI \texttt{text-emb-3-small}     & 0.946          & 0.946          & 0.842 \\
BGE-M3                                & 0.946          & 0.946          & 0.842 \\
\textbf{E5-large}                     & \textbf{0.973} & \textbf{0.973} & \textbf{0.878} \\
ColBERTv2                             & 0.892          & 0.919          & 0.786 \\
\bottomrule
\end{tabular}
\caption{Best per embedder; full grid in \texttt{retrieval\_eval/}.}
\label{tab:retrieval}
\end{table}

E5-large with the reranker wins every paper-level metric, beating the closed-source OpenAI baseline by $+8.6$\,pp MRR; ColBERTv2 trails by $17$\,pp MRR. Multi-hop hit@10 falls from $1.000$ on coarse to $0.667$--$0.833$ on fine: the chunker is the dominant bottleneck.

% TODO (Elie): chunker variation experiment per TA feedback
% ("vary chunk size or overlap, or a semantic/recursive chunker, to close
% some of the multi-hop / open-recall gap"). Add 1 table + 1 plot.

\subsection{Faithfulness verification}
\label{sec:faith}

The verifier extends the proposal's NLI pipeline (DeBERTa-v3-large-mnli, \citealp{he2021debertav3}) with an LLM-judge head. On 24 answerable gold questions answered closed-book by three generators (GPT-4o-mini, Claude-3.5-haiku, DeepSeek-chat), NLI marks $18/24$ ($0.750$) faithful at every provider while the LLM judge marks $0.138$/$0.173$/$0.198$ --- a $55$--$61$\,pp gap matching a known NLI failure mode \citep{min-etal-2023-factscore,gao2023enabling}: NLI alone fails as a faithfulness gate.

% TODO (Andrea + Yusif): human-eval pass per TA feedback ("Your faithfulness
% section flags the 55-61pp NLI-vs-judge disagreement, but doesn't provide
% information about the faithfulness of the judge LLM. Human evaluation
% would be a strong addition to validate your claims here.").
% Target: 60 disagreement items + 20 agreement controls; 2 annotators each;
% report Cohen's kappa human vs (a) NLI, (b) judge.

\subsection{End-to-end RAG vs long-context}
\label{sec:ragas}

Two pipelines run through the four RAGAS metrics \citep{es2024ragas} on an $n{=}8$ stratified subset: chunked RAG (\S\ref{sec:retrieval} retriever, GPT-4o-mini) and a long-context baseline that loads every cited \texttt{document.md} into a 1M-token Gemini~2.5~Flash window, judged by GPT-4o-mini. LC wins every metric but gaps are small (RAG/LC): faithfulness $0.899/0.933$, answer\_relevancy $0.713/0.789$, context\_precision $1.000/1.000$, context\_recall $0.917/1.000$ --- the open recall gap replicates the chunker bottleneck from \S\ref{sec:retrieval}. OpenRouter cost captures show LC at $\sim$10$\times$ the answer-LLM spend (\$0.00383 vs \$0.00038/query), reframing the trade-off as accuracy-per-dollar.

% TODO (Yusif): refresh table with final M3 numbers if n increased beyond 8.

\subsection{Corrective RAG}
\label{sec:crag}

The CRAG router \citep{yan2024corrective} scores top-$k$ retrievals as CORRECT/AMBIGUOUS/INCORRECT and re-queries on AMBIGUOUS. On a mixed 10-question slice ($\tau_h{=}0.813$), CRAG lifts correct-retrieval from 60\% to 80\% (${+}20$~pp), rescuing both ambiguous queries in one refinement round. A 9-cell $(\tau_h, \tau_l)$ sweep hits the same 0.80 rate but spans
0.04--0.40~s/query, with $(\tau_h{=}0.7, \tau_l{=}0.4)$ latency-optimal.

% TODO (Faruk): scale to full 93-question gold set; document INCORRECT-branch
% fallback explicitly (per intermediate feedback) and show one qualitative
% example per branch (CORRECT / AMBIGUOUS / INCORRECT).


\section{Ethical considerations}
\label{sec:ethics}
% TODO (group): per rubric, address when applicable:
%   1. Language adaptation: corpus is English-only; outline how the pipeline
%      could be ported to high-resource (FR, DE) and low-resource (UR, SW)
%      languages (e.g. multilingual embedder swap, MT pivot, native NLI heads).
%   2. Signed-language interaction (rubric prompt) — discuss briefly.
%   3. Beneficiaries vs harms: a faithfulness-gated RAG benefits legal /
%      policy / scientific reviewers who must cite sources; harms include
%      false confidence if the gate fails silently, and downstream effects
%      on IP-litigation workflows.
%   4. Marginalized groups: bias in the source corpus (Anglo-American IP law
%      bias) could under-serve non-Western legal traditions.
%   5. Judge-LLM bias: discuss limitations of using proprietary judges and
%      our planned mitigation (human eval cross-validation).


\section{Conclusion}
\label{sec:conclusion}
% TODO (group): summarize:
%   - Headline result: span-grounded gold benchmark + four-component
%     evaluation isolates retrieval from faithfulness from judging.
%   - Key findings: E5-large coarse+rerank SOTA; chunker is the dominant
%     bottleneck; NLI-vs-judge 55-61pp gap; CRAG +20pp on mixed slice;
%     long-context wins but at 10x cost.
%   - Limitations: small human-eval sample, single English corpus, judge-LLM
%     dependence, CRAG evaluated on 10-question slice (M2) / 93 (M3).
%   - Future work: multi-lingual port, full IAA close-out, scaling CRAG
%     fallback to live web search.


% Entries for the entire Anthology, followed by custom entries
\bibliography{anthology,custom}
\bibliographystyle{acl_natbib}

\appendix

\section{Appendix (optional)}\label{sec:appendix}
% Optional appendix: additional ablation tables, full retrieval grid,
% qualitative CRAG traces, RAGAS per-question breakdowns. Does NOT count
% toward the 4-page limit. Graders are not required to read it.


\end{document}
